\bookchapter{บทที่ 45: การประกันคุณภาพ—ทำให้โค้ดที่ AI เขียนไว้ใจได้}

\chapterintro{โค้ดที่ AI เขียนมีจุดเด่นอย่างหนึ่ง: มันมัก “ดูถูกต้อง”—ตั้งชื่อตัวแปรเป็นระเบียบ คอมเมนต์ครบ ตรรกะชัดเจน แต่ “ดูถูก” ไม่เท่ากับ “ถูกจริง” ปัญหาที่ซ่อนในโค้ด AI มักหายากกว่าโค้ดที่คนเขียนเสียอีก เพราะคุณภาพผิวของมันสูงเกินไป ทำให้คนตายใจ บทนี้สอนวิธีสร้างระบบประกันคุณภาพที่ได้ผลให้โค้ด AI—ด้วยสามแนวป้องกันคือการทดสอบ การรีวิว และการควบคุมเวอร์ชัน เปลี่ยน “ดูถูก” ให้เป็น “ถูกจริง”}

\booksection{ขับเคลื่อนด้วยการทดสอบ: ให้ AI เขียนชุดทดสอบ}

การทดสอบคือแนวป้องกันแรกของการประกันคุณภาพโค้ด AI ข่าวดีคือ: ให้ AI เขียนชุดทดสอบง่ายกว่าและไว้ใจได้กว่าให้มันเขียนโค้ดฟังก์ชัน—เพราะพฤติกรรมที่คาดหวังของการทดสอบชัดเจน โอกาสเกิดโค้ด “ดูถูกแต่จริง ๆ ผิด” จึงน้อยกว่า

ขั้นตอนการเขียนโปรแกรม AI แบบขับเคลื่อนด้วยการทดสอบคือ: ให้ AI เขียนโค้ดฟังก์ชันก่อน พอตรวจฟังก์ชันพื้นฐานผ่าน รีบให้มันเขียนกรณีทดสอบ กรณีทดสอบควรครอบคลุมเส้นทางปกติ กรณีขอบเขต และกรณีผิดปกติ จากนั้นให้ AI แก้โค้ดฟังก์ชันตามฟีดแบ็กจากการทดสอบ ข้อดีของขั้นตอนนี้คือ: พอชุดทดสอบเสร็จ มันคือเครื่องมือตรวจสอบถาวร—การแก้ไขครั้งต่อ ๆ ไปรันชุดทดสอบยืนยันได้เลยว่าไม่แทรกปัญหาเก่ากลับมา

\begin{lstlisting}
ช่วยเขียน unit test ให้ฟังก์ชัน calculate_discount
ชุดทดสอบควรครอบคลุม:
1. กรณีปกติ:
   - สินค้า 100 หยวน ลด 20% -> 80
   - สินค้า 200 หยวน ลด 50% -> 100
2. กรณีขอบเขต:
   - สินค้า 0 หยวน ลด 20% -> 0
   - สินค้า 100 หยวน ไม่ลดเลย -> 100 (ไม่มีส่วนลด)
   - สินค้า 100 หยวน ลด 0% -> 100 (ราคาเต็ม)
3. กรณีผิดปกติ:
   - ราคาติดลบ -> ต้อง throw error
   - อัตราส่วนลด > 1 -> ต้อง throw error
   - ใส่ค่าที่ไม่ใช่ตัวเลข -> ต้อง throw error
ใช้เฟรมเวิร์ก Vitest และให้ไฟล์เทสต์แบบเต็มมา
\end{lstlisting}

สังเกตจุดเด่นของคำสั่งนี้: ระบุกรณีทดสอบสามประเภทชัดเจน (ปกติ ขอบเขต ผิดปกติ) ไม่เปิดช่องให้ AI ใช้ดุลพินิจไปทำ “ทดสอบเลือกบางส่วน” บางครั้ง AI จะเขียนแค่เส้นทางปกติเพื่อ “ส่งงานให้จบไว” การระบุทุกประเภทล่วงหน้าช่วยกันปัญหานี้

\booksection{เทคนิครีวิวโค้ด}

แม้มีชุดทดสอบครอบคลุม คุณก็ยังต้องรีวิวโค้ดที่ AI ผลิตด้วยสายตาคน การรีวิวโค้ดไม่ใช่การอ่านทีละบรรทัด—นั่นเปลืองเวลาและพึ่งทักษะเขียนโปรแกรมของคุณมากเกินไป การรีวิวโค้ด AI ที่ได้ผลควรโฟกัสมิติสำคัญดังนี้:

\begin{center}
\begin{tabularx}{\linewidth}{|X|X|X|}
\hline
\textbf{มิติการรีวิว} & \textbf{ดูอะไร} & \textbf{ปัญหาที่พบบ่อย} \tabularnewline
\hline
ความปลอดภัย & การตรวจสอบอินพุต การจัดการคีย์ การป้องกัน SQL injection & AI มักลืมตรวจอินพุต และเขียนคีย์แปะในโค้ด \tabularnewline
การจัดการข้อผิดพลาด & ข้อยกเว้นถูกจับได้ไหม ข้อความ error เหมาะสมไหม & AI ชอบใช้บล็อก catch เปล่า หรือแค่ console.log \tabularnewline
ประสิทธิภาพ & มีปัญหาประสิทธิภาพเห็นได้ชัดไหม เช่น คิวรีแบบ N+1 ลูปไม่สิ้นสุด & บางครั้ง AI เขียนอัลกอริทึมแบบ O(n²) ออกมา \tabularnewline
การดูแลรักษา & โครงสร้างโค้ดชัดไหม การตั้งชื่อมีความหมายไหม & AI บางครั้งสร้างโค้ดนามธรรมเกินจำเป็น \tabularnewline
การจัดการการพึ่งพา & การพึ่งพาที่เพิ่มเข้ามาจำเป็นไหม เวอร์ชันระบุชัดไหม & AI มักเพิ่มไลบรารีที่ไม่จำเป็น \tabularnewline
\hline
\end{tabularx}
\end{center}

\begin{center}\small\color{gray}มิติสำคัญในการรีวิวโค้ด AI\end{center}

ตอนรีวิวไม่ต้องไล่เช็กเองทีละอย่าง—ให้ AI ช่วยรีวิวรอบแรกได้ คำสั่งประมาณว่า: “ช่วยรีวิวความปลอดภัย การจัดการข้อผิดพลาด และประสิทธิภาพของโค้ดชุดนี้ แล้วชี้ปัญหาที่อาจมี” ความสามารถตรวจตัวเองของ AI ดีขึ้นเรื่อย ๆ แม้แทนการรีวิวโดยคนไม่ได้ทั้งหมด แต่ช่วยกรองปัญหาเห็นได้ชัดส่วนใหญ่ออกได้

\booksection{การควบคุมเวอร์ชัน: บทบาทของ Git ในการเขียนโปรแกรมด้วย AI}

บทบาทของ Git ในการเขียนโปรแกรมด้วย AI สำคัญกว่าในการเขียนโปรแกรมแบบเดิมเสียอีก เพราะการแก้แต่ละครั้งของ AI อาจแทรกความเปลี่ยนแปลงที่คาดไม่ถึง คุณจึงต้องมียาแก้ที่ไว้ใจได้—การควบคุมเวอร์ชันก็คือยาแก้นั้น

แนวปฏิบัติที่ดีที่สุดของ Git ในการเขียนโปรแกรม AI คือ “คอมมิตบ่อย ก้าวเล็กก้าวไว” ทำฟังก์ชันจุดหนึ่งที่ตรวจได้เสร็จ ก็คอมมิตหนึ่งครั้ง วิธีนี้เมื่อการแก้ครั้งใดของ AI มีปัญหา คุณย้อนกลับไปสถานะที่ใช้ได้ก่อนหน้าได้แม่นยำ โดยไม่ต้องล้มทิ้งทั้งหมด

เทคนิคอีกอย่างคือ “แยกสาขา” สำหรับการแก้ด้วย AI ที่ไม่แน่ใจผล (เช่น ปรับใหญ่ทีเดียว) สร้างสาขาใหม่มาทำก่อน พอใจผลก็ merge กลับสาขาหลัก; ไม่พอใจก็ลบสาขาทิ้งตรง ๆ ไม่กระทบสาขาหลัก วิธีนี้ปลอดภัยกว่าการแก้วนไปวนมาบนสาขาหลักมาก

\begin{enumerate}
  \item ก่อนให้ AI แก้ทุกครั้ง ให้มั่นใจว่าพื้นที่ทำงานสะอาด (ไม่มีการเปลี่ยนแปลงที่ยังไม่ได้คอมมิต) วิธีนี้ถ้าผลการแก้ของ AI ไม่พอใจ ใช้ git checkout . ย้อนกลับได้ไว
  \item หลัง AI แก้เสร็จและตรวจผ่านแล้ว คอมมิตทันที ข้อความคอมมิตบอกว่าครั้งนี้แก้อะไร เช่น "feat: เพิ่มฟังก์ชันสมัครสมาชิก"
  \item สำหรับการปรับใหญ่ทีเดียว สร้างสาขาก่อน: git checkout -b refactor/auth-system ทำการแก้บนสาขา ทดสอบผ่านแล้วค่อย merge
  \item ฝึกนิสัยดู diff ทุกครั้งหลัง AI แก้ ใช้ git diff ดูว่าแก้อะไรไปบ้าง สิ่งนี้ไม่เพียงช่วยให้คุณเจอปัญหา แต่ยังช่วยให้คุณเรียนรู้โค้ดด้วย
\end{enumerate}

\booksection{กับดักที่พบบ่อยและคู่มือเลี่ยงกับดัก}

การเขียนโปรแกรมด้วย AI มีกับดักที่เจอซ้ำ ๆ ไม่กี่อัน มือใหม่แทบทุกคนเหยียบเข้าให้ รู้ล่วงหน้าจะช่วยให้คุณไม่ต้องเหนื่อยวนกับมัน

\booksubsection{กับดักที่หนึ่ง: กับดักการวางใจ}

“โค้ดที่ AI เขียนดูมืออาชีพมาก คงไม่มีปัญหาหรอก?”—นี่คือกับดักที่อันตรายที่สุด โค้ดที่ AI สร้างสะอาดระดับผิวมาก—ตั้งชื่อเป็นระเบียบ คอมเมนต์ครบ โครงสร้างชัด—ความ “ดูมืออาชีพที่ผิว” แบบนี้ทำให้คนตายใจ แต่ความสะอาดที่ผิวไม่ได้แปลว่าตรรกะถูก ยาแก้: ตรวจยืนยันด้วยการทดสอบเสมอ ไม่ใช่ตัดสินจาก “ดูถูกไหม”

\booksubsection{กับดักที่สอง: กับดักการบวม}

“เผื่อ AI ช่วยเขียนโค้ดได้ งั้นทำฟีเจอร์ให้ครบขึ้นอีกหน่อยสิ?”—สิ่งนี้ทำให้โค้ดบวม AI เก่งเรื่อง “ให้เกินขอ”: ให้ทำหน้าเข้าสู่ระบบ มันให้ระบบจัดการสมาชิกทั้งชุดมา ฟีเจอร์เกินมาไม่เพียงเพิ่มความซับซ้อน แต่ยังเพิ่มโอกาสเกิดบั๊ก ยาแก้: ยึดหลัก Simplicity First เคร่งครัด บอก AI ชัดว่า “ทำแค่ฟังก์ชันนี้ ห้ามเพิ่มอย่างอื่น”

\booksubsection{กับดักที่สาม: กับดักการเข้าใจ}

“AI คงเข้าใจที่ฉันจะสื่อนะ?”—ไม่แน่ แม้ความสามารถเข้าใจของ AI จะแกร่ง แต่กับความต้องการกำกวม มันยังต้อง “เดา” พอเดาผิด โค้ดที่ผลิตออกมาก็ผิด ที่ยุ่งยากกว่านั้นคือโค้ดที่มันเดาผิดมัก “ดูสมเหตุสมผล” ทำให้คุณเผลอคิดว่ามันเข้าใจ ยาแก้: ให้รายละเอียดเฉพาะเจาะจงพอในคำสั่ง ลดช่องว่างที่ AI ต้อง “เดา”

\booksubsection{กับดักที่สี่: กับดักบริบท}

“ฉันบอกไปแล้วในบทสนทนาก่อนหน้า AI คงยังจำได้?”—มีโอกาสสูงที่มันไม่จำ บทสนทนายาวขึ้นหรือเปิดเซสชันใหม่ AI จะสูญเสียบริบทช่วงต้น ถ้าข้อจำกัดสำคัญข้อใดไม่ได้เขียนลงไฟล์ความจำโปรเจกต์ มันมีโอกาสถูกละเลยในการแก้รอบต่อ ๆ ไป ยาแก้: เขียนข้อจำกัดสำคัญลงไฟล์ความจำโปรเจกต์ อย่าฝากไว้กับความจำของบทสนทนา

\begin{tcbwarn}{จำสี่กับดักใหญ่}
กับดักการวางใจ—อย่าเชื่อ “ดูถูก”; กับดักการบวม—AI ชอบ “ให้เกินขอ”; กับดักการเข้าใจ—ความต้องการกำกวมมันจะ “เดา”; กับดักบริบท—ความจำบทสนทนาจะ “ลืม” ยาแก้ตามลำดับคือ: ตรวจด้วยการทดสอบ, ตัดส่วนเกิน, ให้รายละเอียดเต็มพิกัด, เขียนลงไฟล์ความจำ

\end{tcbwarn}

\begin{tcbpractice}{ลงมือทำ: สร้างสามแนวป้องกันให้โปรเจกต์ของคุณ}
ถ้าคุณเริ่มทำโปรเจกต์ที่หนึ่ง (เว็บเครื่องมือ) แล้ว: (1) ให้ AI เขียนกรณีทดสอบ 3 ประเภทให้ฟังก์ชันหลัก; (2) รีวิวโค้ดตามห้ามิติ (ความปลอดภัย/การจัดการข้อผิดพลาด/ประสิทธิภาพ/การดูแลรักษา/การพึ่งพา); (3) ใช้ git ไล่ตั้งโปรเจกต์และทำคอมมิตแรก ครบสามแนวป้องกันแล้ว โปรเจกต์ของคุณถึงจะเป็นของคุณจริง ๆ

\end{tcbpractice}

\begin{tcbsummary}{สรุปบทนี้}
\begin{itemize}
  \item โค้ด AI "ดูถูก" $\neq$ "ถูกจริง"—คุณภาพผิวสูงเกินไปกลับทำให้ตายใจง่าย
  \item ให้ AI เขียนชุดทดสอบ: ระบุกรณีปกติ/ขอบเขต/ผิดปกติครบสามประเภท กัน "ทดสอบเลือกบางส่วน"
  \item รีวิวห้ามิติ: ความปลอดภัย การจัดการข้อผิดพลาด ประสิทธิภาพ การดูแลรักษา การจัดการการพึ่งพา
  \item Git คือ "ยาแก้": คอมมิตบ่อย ก้าวเล็กก้าวไว แยกสาขา หมั่นดู diff
  \item สี่กับดักใหญ่: วางใจ (ตรวจด้วยการทดสอบ), บวม (ตัดส่วนเกิน), เข้าใจ (ให้รายละเอียดเต็ม), บริบท (เขียนลงไฟล์ความจำ)
\end{itemize}

\end{tcbsummary}
